\subsection{PART A – Review Questions}

\subsubsection{Define the following terms in your own words}

\noindent\textbf{Project}: 


A temporary endeavor with a defined beginning and end, undertaken to create a unique product, service, or result.


\noindent\textbf{Project Management Methodology}: 

A structured approach (e.g., Waterfall, Agile, DevOps) that provides processes, practices, and guidelines for planning, executing, and delivering projects.

\noindent\textbf{IT Lifecycle}: 


The stages an IT system goes through, from planning and development to deployment, monitoring, and retirement.

\noindent\textbf{Continuous IT Lifecycle}: 


An iterative model where IT development, deployment, and monitoring occur continuously, with ongoing updates and improvements.

\noindent\textbf{Enterprise Architecture (EA)}: 


A framework for aligning IT strategy, processes, information, and technology with business goals.


\subsubsection{List three common reasons projects fail and explain why}


\begin{itemize}
    \item \textbf{Poor scope definition}: Leads to misunderstandings and scope creep.
    \item \textbf{Unrealistic timelines/budgets}: Causes delays and resource strain.
    \item \textbf{Weak stakeholder engagement}: Results in lack of buy-in and misaligned outcomes
\end{itemize}

\newpage

\subsection{PART B – Group Discussion}

\subsubsection{Compare Waterfall, Agile, and DevOps}


\begin{table}[h]
\centering
\begin{tabular}{|cl|ll|ll|l|}
\hline
\multicolumn{2}{|c|}{\textbf{Aspect}}                                                            & \multicolumn{2}{c|}{\textbf{Waterfall}}                                                                                       & \multicolumn{2}{c|}{\textbf{Agile}}                                                                 & \multicolumn{1}{c|}{\textbf{DevOps}}                                           \\ \hline
\multicolumn{2}{|c|}{\textbf{Flexibility}}                                                       & \multicolumn{2}{l|}{\begin{tabular}[c]{@{}l@{}}Low-rigid, hard to\\ change\end{tabular}}                                      & \multicolumn{2}{l|}{\begin{tabular}[c]{@{}l@{}}High-adapts to evolving\\ requirements\end{tabular}} & \begin{tabular}[c]{@{}l@{}}High-continuous\\ improvements\end{tabular}         \\ \hline
\multicolumn{2}{|c|}{\textbf{Delivery speed}}                                                    & \multicolumn{2}{l|}{\begin{tabular}[c]{@{}l@{}}Slow-end of project\\ delivery\end{tabular}}                                   & \multicolumn{2}{l|}{\begin{tabular}[c]{@{}l@{}}Faster-incremental\\ delivery\end{tabular}}          & \begin{tabular}[c]{@{}l@{}}Fastest-continuous\\ deployment\end{tabular}        \\ \hline
\multicolumn{2}{|c|}{\textbf{\begin{tabular}[c]{@{}c@{}}Stakeholder\\ Involvement\end{tabular}}} & \multicolumn{2}{l|}{\begin{tabular}[c]{@{}l@{}}Low-mainly at start\\ and end\end{tabular}}                                    & \multicolumn{2}{l|}{\begin{tabular}[c]{@{}l@{}}High-continuous\\ collaboration\end{tabular}}        & \begin{tabular}[c]{@{}l@{}}High-ongoing feedback \\ \& monitoring\end{tabular} \\ \hline
\multicolumn{2}{|c|}{\textbf{\begin{tabular}[c]{@{}c@{}}Best-fit\\ Examples\end{tabular}}}       & \multicolumn{2}{l|}{\begin{tabular}[c]{@{}l@{}}Compliance/regulatory\\ projects (e.g., NASA\\ Shuttle software)\end{tabular}} & \multicolumn{2}{l|}{\begin{tabular}[c]{@{}l@{}}Evolving projects like\\ SaaS apps\end{tabular}}     & \begin{tabular}[c]{@{}l@{}}Large-scale platforms\\ (e.g., Amazon)\end{tabular} \\ \hline
\end{tabular}
\end{table}


\subsubsection{Why is the shift from a linear IT lifecycle to a continuous IT lifecycle important in today’s digital environment?}


\begin{itemize}
    \item \textbf{Digital transformation}: Organizations must respond faster to market and technology changes.
    \item \textbf{Customer expectations}: Users demand quick updates, bug fixes, and new features.
    \item \textbf{Cloud and automation}: Enable faster, safer deployments.
    \item \textbf{Agile \& DevOps practices}: Promote collaboration, speed, and iterative delivery
\end{itemize}


\subsection{PART C – Case Study Discussion}

\indent 

\noindent\textbf{Case Scenario}:

HealthLink Systems (fictional healthcare software provider) faced frequent delays in delivering updates due to siloed IT and business teams. Customers complained about long waits for bug fixes and feature requests. HealthLink shifted from a Waterfall model to a Continuous IT Lifecycle supported by TOGAF-based Enterprise Architecture, achieving faster releases, reduced costs, and improved customer satisfaction.

\bigskip

\noindent\textbf{Issues faced before the shift}:

\begin{itemize}
    \item Frequent delays in updates due to siloed IT/business teams.
    \item Customer dissatisfaction with long wait times for fixes.
    \item Inefficient delivery under the Waterfall model.
\end{itemize}


\bigskip

\noindent\textbf{Why was moving to a continuous IT lifecycle more suitable?}


\begin{itemize}
    \item Allowed incremental updates, faster bug fixes, and quicker feature delivery.
    \item Reduced costs through automation and streamlined processes.
    \item Improved collaboration across IT and business teams
\end{itemize}


\bigskip

\noindent\textbf{Would Agile alone have solved the problem, or was DevOps also necessary?}:


\begin{itemize}
    \item Agile would improve responsiveness but DevOps was necessary for automation, continuous integration, and faster deployment.
    \item The combination ensured both adaptability and operational efficiency.
\end{itemize}


\bigskip

\noindent\textbf{How did Enterprise Architecture (EA) contribute to long-term success?}:

\begin{itemize}
    \item Provided a structured framework (TOGAF) for aligning IT with business goals.
    \item Ensured integration of applications, data, and technology architecture across the organization.
    \item Supported scalability and ongoing innovation
\end{itemize}

\bigskip

\noindent\textbf{How could EA support adding AI-based predictive analytics?}:

\begin{itemize}
    \item Ensures smooth integration of AI tools with existing hospital systems.
    \item Provides standards for data architecture and interoperability.
    \item Reduces risks by aligning new technology with strategic goals.
\end{itemize}